\documentclass[a4paper,11pt]{article}
% Подключаем все нужные пакеты
%%% Локализация и шрифты
\usepackage[utf8]{inputenc}          % Кодировка UTF-8
\usepackage[T2A]{fontenc}            % Русская кодировка
\usepackage[english,russian]{babel} % Локализация
\usepackage{cmap}                    % Поддержка поиска в PDF
\usepackage{tikz}
\usepackage{xstring}
\usepackage[left=20mm, top=10mm, right=20mm, bottom=10mm, nohead, nofoot]{geometry}
\renewcommand{\baselinestretch}{1.0}% Межстрочный интервал
\usepackage{ragged2e}               % Выравнивание текста
\usepackage{microtype}              % Улучшение качества вывода текста
\usepackage{multicol}
\usepackage{tcolorbox}
\tcbuselibrary{listingsutf8, breakable, skins}
\justifying
\sloppy
\tolerance=500
\hyphenpenalty=10000
\emergencystretch=3em

\usepackage[normalem]{ulem}
\pagestyle{empty} \textwidth=19.0cm \oddsidemargin=-1.3cm
\textheight=26cm \topmargin=-3.0cm
\usepackage{tcolorbox}
\tcbuselibrary{listingsutf8}
\usepackage{pgfplots}
\renewcommand{\rmdefault}{cmss}
\usepgfplotslibrary{patchplots}
\pgfplotsset{compat=1.18}
\usetikzlibrary{fillbetween}
% Создаём новый счётчик задач
\newcounter{mainth}
\newcounter{subth}[mainth]

\newcommand{\thsub}{\arabic{mainth}.\arabic{subth}}

\newcounter{mainpr}
\newcounter{subpr}[mainpr]

\newcommand{\propsub}{\arabic{mainpr}.\arabic{subpr}}

\newcounter{mainex}
\newcounter{subex}[mainex]

\newcommand{\exsub}{\arabic{mainex}.\arabic{subex}}

\renewcommand{\familydefault}{\ttdefault}

% Определяем стили блоков
\usepackage{xcolor}

% --- Однострочное создание блоков \problem, \solution, \answer ---
% (работает без \begin ... \end)
\definecolor{myrose}{RGB}{255,182,193} % Розовый
\definecolor{mybrown}{RGB}{210,180,140} % Светло-коричневый
\definecolor{mycyan}{RGB}{135,206,235}  % Голубой
\definecolor{myorange}{RGB}{255,165,0}   % Оранжевый
\definecolor{mypurple}{RGB}{216,191,216} % Светло-фиолетовый
% Требуется пакет environ:
\newcommand{\thletter}{}

\newcommand{\thtitle}{
  \IfStrEq{\thletter}{}%  % Если taskletter пустая
    {\ttfamily}%   % Выводим только номер
    {\ttfamily (\thletter)}% Иначе номер + буква
}

% --- Определяем сокращённые версии окружений ---

\newtcolorbox{defBox}{
  breakable,
  enhanced,
  colback=violet!10!white,
  colframe=mypurple,
  fonttitle=\ttfamily,
  title={Определение}
}

\newcommand{\defin}[1]{%
  \begin{defBox}%
    #1%
  \end{defBox}%
}

\newtcolorbox{theoremBox}{
  breakable,
  enhanced,
  colback=red!10!white,
  colframe=myrose,
  fonttitle=\ttfamily,
  title={Теорема~\thsub \thtitle}
}

\newcommand{\theorem}[1]{%
  \refstepcounter{subth}%
  \begin{theoremBox}%
    #1%
  \end{theoremBox}%
}

\newtcolorbox{proposBox}{
  breakable,
  enhanced,
  colback=brown!10!white,
  colframe=mybrown,
  fonttitle=\ttfamily,
  title={Утверждение~\propsub}
}

\newcommand{\propos}[1]{%
  \refstepcounter{subpr}%
  \begin{proposBox}%
    #1%
  \end{proposBox}%
}

\newtcolorbox{proofBox}{
  breakable,
  enhanced,
  colback=cyan!10!white,
  colframe=mycyan,
  fonttitle=\ttfamily,
  title={Доказательство}
}

\newcommand{\prof}[1]{
  \begin{proofBox}
    #1
  \end{proofBox}
}

\newtcolorbox{exampleBox}{
  breakable,
  enhanced,
  colback=orange!10!white,
  colframe=myorange,
  fonttitle=\ttfamily,
  title={Пример~\exsub}
}

\newcommand{\example}[1]{%
  \refstepcounter{subex}%
  \begin{exampleBox}%
    #1%
  \end{exampleBox}%
}

% Команда для векторов
\newcommand{\vect}[1]{\mathbf{#1}}
\usepackage[upint]{stix2}

%%% Колонтитулы
\usepackage{fancyhdr}
\pagestyle{fancy}
\renewcommand{\headrulewidth}{0.3mm} % Толщина линии в колонтитулах
% ПРОВЕРЬТЕ ИМЯ ФАЙЛА! Должно быть 'profimatika' или 'profimatica'

\rhead{Теория вероятностей и математическая статистика}
\lhead{Благий Александр}
\headsep=25pt
\footskip=15mm
\textheight=700pt
\headheight=35pt % Возможно, понадобится увеличить до ~30pt. Проверьте log-файл.
\parindent=0mm
%%% Математика
\usepackage{amsmath, amsfonts, amssymb, amsthm, mathtools} % Базовые математические пакеты
\usepackage{icomma}   
\usepackage{euscript}               % Шрифт Евклид
\usepackage{mathrsfs}               % Красивый матшрифт
\usepackage{cancel}                 % Зачёркивание в формулах

% Замены часто употребляемых шрифтов
\newcommand{\mb}[1]{\mathbb{#1}}
\newcommand{\mc}[1]{\mathcal{#1}}
\newcommand{\ms}[1]{\mathscr{#1}}